% unit U1
\section{Introduction}\label{sec:intro}
\mathplain

Scala 3 can check at compile time which capabilities a value may use \cite{scala3cc,boruchgruszecki2023capturing}. A capability is a value that gives access to an effect, such as a file handle or a logger. The checker writes capabilities into types, and it rejects a program whose types do not account for the capabilities it uses. This report explains why that verdict can be trusted, for capture checking over Dependent Object Types (DOT), by the route that \Rone took for type safety.

\paragraph{Capture checking.}
The page on capture checking basics in the Scala 3 reference \cite{scala3cc} opens with a method that lends a log file to an operation and closes the file when the operation returns. We quote the reference at commit \texttt{4dae25087d} throughout, by file and line. This is the checked version, \texttt{basics.md}, lines 56 to 60:
\begin{lstlisting}[language=Scala]
def usingLogFile[T](op: FileOutputStream^ => T): T =
  val logFile = FileOutputStream("log")
  val result = op(logFile)
  logFile.close()
  result
\end{lstlisting}
The \lstinline|^| after the parameter type makes the file a capability whose uses the checker tracks. An operation that returns a closure over the file is rejected, with this message (lines 65 to 70):
\begin{lstlisting}[basicstyle=\ttfamily\footnotesize,columns=fixed,basewidth=0.52em]
   |  val later = usingLogFile { f => () => f.write(0) }
   |                           ^^^^^^^^^^^^^^^^^^^^^^^^^
   |  Found:    (f: java.io.FileOutputStream^'s1) ->'s2 () ->{f} Unit
   |  Required: java.io.FileOutputStream^ => () ->'s3 Unit
   |
   |  Note that capability f cannot be included in outer capture set 's3.
\end{lstlisting}
The closure \lstinline|() => f.write(0)| has the type \lstinline|() ->{f} Unit|, which says that it may use \lstinline|f|. Called after \lstinline|usingLogFile| has returned, it would write to a closed file. In the message, \lstinline|'s1| to \lstinline|'s3| are capture sets that the checker infers. The set \lstinline|'s3| belongs to the result type \lstinline|T|, which is fixed where \lstinline|usingLogFile| is called, outside the scope of \lstinline|f|. So \lstinline|'s3| cannot contain \lstinline|f|, and the program is rejected before it runs. The verdict comes from types alone, so it can be trusted only if the types are right about runs. A value whose type lists a set of capabilities must use no others when the program runs. This report proves a statement of that kind. The escaped closure uses only \lstinline|f|, and its type lists \lstinline|f|. So the statement does not rule out the escape, which is a matter of scope. Section~\ref{sec:levels} treats the escape and says how much of it is proved.

\paragraph{What is proved.}
This report proves effect safety, Theorem~\ref{thm:mini-effect} in the miniature and \lean{effect\_safety} in each of the three Lean trees. Let $\kappa$ be a capture variable of the platform, which stands for a primitive capability such as a console. If $\kappa$ is not a root of the use set of a typed state, then no state of its run inspects a variable $x$ with $\kappa$ among the roots of $\cset{x}$. The use set of a state is what the state may still use, read off its syntax. The roots of a capture set are the capture variables it reaches when each variable in it is replaced by the capture set of its type. Later sections extend this definition. A state inspects the variable whose stored value its next step reads. Section~\ref{sec:cc} explains the three words, Section~\ref{sec:mini} defines them, and Theorem~\ref{thm:mini-effect} states the theorem with all its hypotheses. In the running example of Section~\ref{sec:mini}, a program calls a closure that writes to a log. It drops the closure that this call returns, which would write to a console. The roots of its use set include the capture variable of the log and not that of the console, so its run never calls the console (Example~\ref{ex:mini-prog2}). Effect safety follows from capture prediction, which says that the roots of the use set never grow along a run (Section~\ref{sec:mini-run}). Section~\ref{sec:source-thms} carries both theorems to DOT programs over a platform, through a target state with the same erasure as the source state.

\paragraph{The route.}
Report~1 proved DOT type safe by a route that compilers take. A typing derivation is translated into a term of \FCdot, in which every use of subtyping is an explicit proof term, called evidence, that erases. Type safety of \FCdot rests on one theorem about closed evidence over a typed store (Section~\ref{R1-sec:idea} of report~1). Safety of DOT follows by running the source program and its translation side by side. This report adds capture sets to both calculi and keeps the route. Subcapturing becomes evidence too. Every binder declares the use set of its body and carries evidence that the body stays within it, so the use set of a term is computed from its syntax. One more theorem about closed evidence says that over a typed store, capture evidence that $C$ is below $D$ puts every root of $C$ among the roots of $D$ (Theorem~\ref{thm:target-canon}). A use-set half of preservation then reads the evidence that each step needs as an inclusion of roots. The source enters by translation. A DOT derivation becomes an \FCdot term together with evidence that its use set fits the declared one, and a source run is matched by a target run with the same erasure. The proof is syntactic, by preservation and progress, and uses no logical relation.

\paragraph{Three trees.}
The development has three Lean trees. \treeK is the main tree. \treeCC adds the compiler's scopes and levels, and \treeCL adds capability classifiers (Section~\ref{sec:mech-trees}).

\paragraph{What is mechanized.}
Every theorem stated in Sections~\ref{sec:target} to~\ref{sec:frontends} is a Lean declaration. Its head, or each of its items, names the declaration, and its head names the tree. The miniature of Section~\ref{sec:mini}, a fragment of the \treeK target, is not mechanized on its own (Section~\ref{sec:mech-mini}). Every example is marked in one of three ways. It is decided when the kernel closes it by \lean{decide} or \lean{rfl}, or by running a checker whose soundness theorem turns the verdict into a derivation. It is built by hand when its Lean proof is written out. It is on paper when the Lean does not state it. The three trees contain no \lean{sorry}, no \lean{axiom} declaration and no \lean{native\_decide}, and the main theorems depend on propositional extensionality and quotient soundness only.

\paragraph{Reading paths.}
Section~\ref{sec:cc} introduces capture checking for a reader who has never seen it, and Section~\ref{sec:problem} shows why the textbook proof needs an inversion that capture members break. Section~\ref{sec:mini} shows the shape of the whole proof on a fragment without capture members, with one running example. Section~\ref{sec:recap} recalls the parts of report~1 that the full calculus uses. Sections~\ref{sec:target} to~\ref{sec:classifiers} scale the proof up to the full \treeK target, the source and its translation, scopes and levels in \treeCC, and classifiers in \treeCL. Section~\ref{sec:frontends} describes the front ends that turn programs into derivations. Section~\ref{sec:history} records the designs that did not work, Section~\ref{sec:mech} is a tour of the Lean, and Sections~\ref{sec:related} and~\ref{sec:conclusion} discuss related work and what is left out. A reader who wants the idea and not the details can read Sections~\ref{sec:cc} to~\ref{sec:mini} and the first paragraph of Section~\ref{sec:conclusion}. A reader who knows Scala's checker can skim Section~\ref{sec:cc} up to the roots of a capture set, and Section~\ref{sec:cc-later}. A Lean reader can start at Section~\ref{sec:mech}.

\paragraph{Conventions.}
As in report~1, source syntax and judgments are set in \src{blue}, target syntax and judgments in \tgt{red}, and the untyped runtime in black. Lean names are set in \lean{typewriter}. The Lean scopes syntax by de Bruijn indices, and the report names every binder instead. Target rules are named in small caps with their judgment as a prefix, such as \rn{C-Var} for capture evidence and \rn{E-App} for a step. Source rules keep the names of report~1, and source subcapturing rules take the prefix \rn{Sc-}.
